\newpage
\begin{center}
\large\textbf{Минобрнауки России}

\large\textbf{Юго-Западный государственный университет}
\vskip 1em
\normalsize{Кафедра программной инженерии}
\vskip 1em
\ifВКР{
        \begin{flushright}
        \begin{tabular}{p{.4\textwidth}}
        \centrow УТВЕРЖДАЮ: \\
        \centrow Заведующий кафедрой \\
        \hrulefill \\
        \setarstrut{\footnotesize}
        \centrow\footnotesize{(подпись, инициалы, фамилия)}\\
        \restorearstrut
        «\underline{\hspace{1cm}}»
        \underline{\hspace{3cm}}
        20\underline{\hspace{1cm}} г.\\
        \end{tabular}
        \end{flushright}
        }\fi
\end{center}
\vspace{1em}
  \begin{center}
  \large
\ifВКР{
ЗАДАНИЕ НА ВЫПУСКНУЮ КВАЛИФИКАЦИОННУЮ РАБОТУ
  ПО ПРОГРАММЕ БАКАЛАВРИАТА}
  \else
ЗАДАНИЕ НА КУРСОВУЮ РАБОТУ (ПРОЕКТ)
\fi
\normalsize
  \end{center}
\vspace{1em}
{\parindent0pt
  Студента \АвторРод, шифр\ \Шифр, группа \Группа
  
1. Тема «\Тема\ \ТемаВтораяСтрока»
\ifВКР{
утверждена приказом ректора ЮЗГУ от \ДатаПриказа\ № \НомерПриказа
}\fi.

2. Срок предоставления работы к защите \СрокПредоставления

3. Исходные данные для создания программной системы:

3.1. Перечень решаемых задач:}

\renewcommand\labelenumi{\theenumi)}

\begin{enumerate}
\item Провести анализ предметной области онлайн-бронирования услуг.
\item Разработать концептуальную модель платформы онлайн-бронирования для частного мастера и клиента.
\item Спроектировать программную систему, базу данных, пользовательский интерфейс и интеграционные взаимодействия.
\item Реализовать и протестировать платформу онлайн-бронирования услуг средствами web-технологий.
\end{enumerate}

{\parindent0pt
  3.2. Входные данные и требуемые результаты для программы:}

\begin{enumerate}
\item Входными данными для программной системы являются: регистрационные данные пользователей; данные профиля мастера и клиента; параметры рабочей доступности мастера; пригласительные ссылки; сведения о выбранных временных интервалах; данные внешнего календаря; настройки каналов уведомлений.
\item Выходными данными для программной системы являются: список доступных слотов для записи; созданные, измененные и отмененные бронирования; сведения о клиентах мастера; статусы приглашений; события синхронизации с внешним календарем; уведомления о событиях записи; эксплуатационные метрики и журналы работы системы.
\end{enumerate}

{\parindent0pt

  4. Содержание работы (по разделам):
  
  4.1. Введение.
  
  4.2. Анализ предметной области.
  
4.3. Техническое задание: основание для разработки, назначение разработки,
требования к программной системе, требования к оформлению документации.

4.4. Технический проект: общие сведения о программной системе, проект
данных программной системы, проектирование архитектуры программной системы, проектирование пользовательского интерфейса и интеграционных взаимодействий.

4.5. Рабочий проект: спецификация компонентов, классов и маршрутов программной системы, тестирование программной системы, сборка компонентов программной системы.

4.6. Заключение.

4.7. Список использованных источников.

5. Перечень графического материала:

Лист 1. Сведения о ВКРБ

Лист 2. Цель и задачи разработки

Лист 3. Диаграмма прецедентов для мастера

Лист 4. Диаграмма прецедентов для клиента

Лист 5. Диаграмма прецедентов для администратора

Лист 6. Общая архитектура системы

Лист 7. Модель данных

Лист 8. Сценарий бронирования

Лист 9. Реляционная модель базы данных

Лист 10. Заключение

\vskip 2em
\begin{tabular}{p{6.8cm}C{3.8cm}C{4.8cm}}
Руководитель \ifВКР{ВКР}\else работы (проекта) \fi & \lhrulefill{\fill} & \fillcenter\Руководитель\\
\setarstrut{\footnotesize}
& \footnotesize{(подпись, дата)} & \footnotesize{(инициалы, фамилия)}\\
\restorearstrut
Задание принял к исполнению & \lhrulefill{\fill} & \fillcenter\Автор\\
\setarstrut{\footnotesize}
& \footnotesize{(подпись, дата)} & \footnotesize{(инициалы, фамилия)}\\
\restorearstrut
\end{tabular}
}

\renewcommand\labelenumi{\theenumi.}
